\documentclass[11pt]{article}
\usepackage{fontspec}
\usepackage[french,english]{babel}
\usepackage{amsmath}
\usepackage{amssymb}
\usepackage{amsthm}
\usepackage{mathrsfs}
\usepackage{tikz-cd}
\usepackage{enumitem}

\newcommand{\printedpage}[1]{\marginpar{\footnotesize p.~#1}}
\newcommand{\msp}[1]{\marginpar{\footnotesize [#1]}}
\newcommand{\Fl}{\mathrm{Fl}}
\newcommand{\Cof}{\mathrm{Cof}}
\newcommand{\deq}{\overset{\text{def}}{=}}
\newcommand{\univ}{\mathrm{univ}}

\begin{document}

\printedpage{24}
% head: A. Grothendieck : Dérivateurs, Ch. XV, Édition des Universités de Montpellier II et Paris VII
% foot: Version 5 novembre 2010

\msp{page 35}
Indeed, by corollary 2 of proposition 4.2, the hypothesis on $TC$ implies that $TC = \widetilde{TC}$, so since $TC \subset C{\cap}W$, the hypothesis $C{\cap}W \subset \widetilde{TC}$ $(= TC)$ is equivalent to $C{\cap}W = TC$.

\textbf{Corollary 2.} \textit{Suppose $C$ and $W$ stable under direct factors. Then the equivalent conditions of proposition} 4.3 \textit{are equivalent to}
\[
C\cap W \;=\; \widetilde{TC}\ . \tag*{(4.3.5)}
\]
\textit{Dually, if $F$, $W$ are stable under direct factors, the equivalent conditions of proposition} 4.3 \textit{are equivalent to}
\[
F\cap W \;=\; \overline{TF}\ . \tag*{(4.3.5$'$)}
\]

Indeed, if $C$ and $W$ are stable under direct factors, the same is true of $C\cap W$, and hence if $C\cap W\subset\widetilde{TC}$, since $C\cap W$ contains $TC$, we conclude $\widetilde{C\cap W}=\widetilde{TC}$. But since $C\cap W\longleftrightarrow TC$ and $C\cap W$ is stable under direct factors, $C\cap W$ is left closed by corollary 2 of proposition 4.2, \textit{i.e.} $C\cap W=\widetilde{C\cap W}$, whence $C\cap W=\widetilde{TC}$, \quad \textit{q.e.d.}

\textbf{Corollary 3.} \textit{Suppose the equivalent conditions} (a) \textit{to} (c) \textit{of} 4.3 \textit{satisfied, and moreover} $W$ stable under direct factors\footnote[12]{In fact, it suffices to suppose that we have $W\supset\widetilde{TC}$, $\overline{TF}$, \textit{i.e.} that every direct factor of an element of $W$ which is either in $C$, or in $F$, is still in $W$. Then $(\widetilde{C},\widetilde{TC},\bar{F},\overline{TF})$ define a closed Quillen structure, and we know that this implies that in fact, $W$ too is stable under direct factors.}. \textit{Then the preliminary conditions} (4.3.1) \textit{and} (4.3.2) \textit{remain satisfied when one replaces $(C,TF)$ and $(TC,F)$ by their closures $(C',TF')=(\widetilde{C},\overline{TF})$ and $(TC',F')=(\widetilde{TC},\bar{F})$, and the equivalent conditions of proposition} 4.3 \textit{are still satisfied for these couples, in the strong forms}
\[
\left\{
\begin{aligned}
C'\cap W &\;=\; TC' && (=\widetilde{TC'}\,)\\
F'\cap W &\;=\; TF' && (=\overline{TF'}\,)\\
W &\;=\; TF'\circ TC'\ .
\end{aligned}
\right.
\tag*{(4.3.6)}
\]

\msp{page 36}
\textit{The triple $(W,C',F')$ is a closed Quillen triple $($corresponding to a ``closed model category'' of Quillen, at least if one assumes $\mathcal{M}$ stable under finite $\varinjlim$ and finite $\varprojlim$, and that $W$ is mildly saturated$)$.}\footnote[13]{This gives the\par
\textbf{Corollary 4.} \textit{Let $W$ be a mildly saturated localizer, and $(C,TF)$, $(TC,F)$ two Quillen precouples such that $TC\subset C\cap W$, $TF\subset F\cap W$. Then the following conditions are equivalent.}
\begin{enumerate}[leftmargin=*]
\item[a)] $TC=C\cap W$.
\item[b)] $TF=F\cap W$.
\item[c)] $W=TF\circ TC$.
\item[d)] $(W,C,F)$ \textit{is a closed Quillen triple.}
\end{enumerate}}

\textsc{Proof.} The relations (4.3.1) imply
\[
\widetilde{TC}\;\subset\;\widetilde{C}\qquad\text{and}\qquad\overline{TF}\;\subset\;\bar{F}\ ,
\]
\printedpage{25}
% head: A. Grothendieck : Dérivateurs, Ch. XV, Édition des Universités de Montpellier II et Paris VII
% foot: Version 5 novembre 2010
\textit{i.e.}
\[
TC'\;\subset\;C'\ ,\qquad TF'\;\subset\;F'\ ,
\]
\textit{i.e.} (4.3.1) for the new Quillen couples. The relation 4.3.2, $W\supset TC,TF$, implies
\[
W\;\supset\;\widetilde{TC}\ ,\ \overline{TF}\ ,\qquad \textit{i.e.}\qquad W\;\supset\;TC'\ ,\ TF'\ ,
\]
by virtue of corollary 2 to proposition 4.2 and the fact that $W$ is stable under direct factors; so we have (4.3.2) for the new couples. Since by hypothesis we have
\[
W\;=\;\overline{TF}\circ\widetilde{TC}\ ,
\]
this reads
\[
W\;=\;TF'\circ TC'\ ,
\]
so the new couples satisfy the same conditions, which this time, the couples being closed, take the strong form (4.3.6).

Let us prove that the triple $(W,C',F')$ satisfies Quillen's axioms for a closed Quillen triple. We do have, by (4.3.6),
\[
\left\{
\begin{aligned}
C' &\;=\;(TF')^*\;=\;(F'\cap W)^*\ , &\qquad F' &\;=\;(TC')_*\;=\;(C'\cap W)_*\ ,\\
W &\;=\;TF'\circ TC'\;=\;C'_*\circ F'^*\ ,
\end{aligned}
\right.
\]
which corresponds to the definition of Quillen's \textit{closed} model categories: Quillen's axioms M\,1, M\,2, M\,3, M\,4 are trivially satisfied (M\,3 and M\,4 because
\msp{page 37}
$C$ and $F$, $TC$ and $TF$ are left closed resp. right closed). Axiom M\,5 is the mild saturation of $W$.

\textbf{4.4.} Here is the idea for constructing more or less ``all'' the closed Quillen structures $(C,TC,F,TF)$ associated to a given localizer $W\subset\mathcal{M}$. We suppose that $\mathcal{M}$ satisfies the conditions of the preceding paragraph, concerning the existence of filtering $\varinjlim$, of a cardinal filtration, accessibility in $\mathcal{M}$ (that of the elements of $\mathrm{Filt}^{\pi_0}$). We start from a small subset
\[
C_0\;\subset\;\Fl(\mathcal{M})
\]
such that
\[
C_{0*}\;\subset\;W\ , \tag*{(4.4.1)}
\]
and we set
\[
\left\{
\begin{aligned}
C &\;=\;\widetilde{C_0}\ , &\qquad TC &\;=\;C\cap W\ ,\\
TF &\;=\;C_*\;=\;C_{0*}\ , &\qquad F &\;=\;(TC)_*\ .
\end{aligned}
\right.
\tag*{(4.4.2)}
\]
There is only one problem, since we already know that $(C,TF)$ is a Quillen couple by the factorization theorem in §1 (p. 10): we need the factorization for the couple $(TC,F\;=\;(TC)_*)$. At that point, the preliminary conditions (4.3.1), (4.3.2) of proposition 4.3 are satisfied:
\[
TC\;\subset C\ ,\qquad TF\;\subset F\ ,\qquad W\;\supset\;TC\ ,\ TF\ ,
\]
as well as condition $(a)$, $C\cap W\subset\widetilde{TC}$, since $C\cap W=TC$. So we find, by the equivalence of conditions $(a)$, $(a')$, $(c)$, and since $TC$ and $TF$ are closed,
\[
TF\;=\;F\cap W\ ,\qquad W\;=\;TF\circ TC\ ,
\]
\printedpage{26}
% head: A. Grothendieck : Dérivateurs, Ch. XV, Édition des Universités de Montpellier II et Paris VII
% foot: Version 5 novembre 2010
\msp{page 38}
on the other hand the hypothesis on $W$ implies that $W$ is stable under direct factors, $C$ is too, being left closed, so $TC=C{\cap}W$ also, so by proposition 4.1, corollary 2, we have $TC\;=\;\widetilde{TC}$, $TC$ is left closed. So we win, by the last assertion of corollary 3, proposition 4.3 (p. 35).

In fact, independently of the particular conditions on $\mathcal{M}$, $W$, we have just proved this:

\textbf{Lemma 4.4.1} (corollary of proposition 4.3.). \textit{Let $\mathcal{M}$ be a category, $W$ a localizer in $\mathcal{M}$ mildly saturated and stable under direct factors. Then the closed Quillen structures $(C,F)$ associated to $W$ are obtained as follows: one takes a subset $C_0$ of} $\Fl(\mathcal{M})$ \textit{such that}
\[
TF\;\deq\;C_{0*}\;\subset\;W\ ,
\]
\textit{one sets}
\[
C\;=\;\widetilde{C_0}\;=\;TF^*\ ,\qquad TC\;=\;C\cap W\ ,\qquad F\;=\;(TC)_*\ ,
\]
\textit{and one must manage to get $(C,TF)$ and $(TC,F)$ to satisfy Quillen's factorization condition,} i.e. \textit{be Quillen precouples.} (In fact, $C$, $TF$ and $F$ being closed by construction, and $TC=C\cap W$ being stable under direct factors ($C$ and $W$ being so), they will even be Quillen \textit{couples}.)

\msp{page 39}
The additional hypotheses made earlier on $\mathcal{M}$ served to ensure that if I take $C_0$ small, then the factorization works for $(C,TF)$ automatically. There remains the problem for $TC$, $F$.

Let us note that if it works, then $TC=\widetilde{TC}=F^*$, so $TC$ is left closed, and \textit{a fortiori} satisfies the corresponding stability conditions:
\[
\left\{
\begin{tabular}{rp{0.6\textwidth}}
\textit{a}) & Stable under cobase change, under composition, under direct factors (these last two are ensured automatically, since $C$, $W$ are stable).\\[1ex]
\textit{b}) & Stable under strict ordinal $\varinjlim$.
\end{tabular}
\right.
\tag*{(4.4.3)}
\]

The stability hypothesis on $W$ serves (among other things) to guarantee stability $(b)$ for $TC$ (in fact, for this alone it would suffice that $W$ be stable under strict ordinal $\varinjlim$). There remains the stability of $TC=C{\cap}W$ under cobase change, which is problematic since $W$ is in general \textit{not} stable under cobase change. It is to ensure this stability that I find myself practically compelled here to impose
\[
C_0\;\subset\;\Cof_W\ , \tag*{(4.4.4)}
\]
I say that with the stability hypothesis made on $W$, this implies
\[
C\;\deq\;\widetilde{C_0}\subset\Cof_W\ . \tag*{(4.4.5)}
\]
\printedpage{27}
% head: A. Grothendieck : Dérivateurs, Ch. XV, Édition des Universités de Montpellier II et Paris VII
% foot: Version 5 novembre 2010
\msp{page 40}
Indeed, we have:

\textbf{Proposition 4.4.2.} \textit{Let $\mathcal{M}$ be stable under small filtering direct $\varinjlim$, and $W\subset$} $\Fl(\mathcal{M})$ \textit{a localizer $($mildly saturated$)$} stable under the said limits. \textit{Then} $\Cof_W$ \textit{satisfies the stability conditions} (4.4.3) (``\textit{left Quillen stability}''), \textit{and is also stable under arbitrary filtering $\varinjlim$.}

Indeed, condition $(a)$ is satisfied without any additional condition on $W$. On the other hand, it is immediate that this condition with respect to filtering $\varinjlim$ implies that $\Cof_W$ \textit{is stable under filtering direct limits} in $\underline{\Fl}(\mathcal{M})$.

\textbf{Corollary 4.4.3.} \textit{Under the conditions of proposition} 4.4.2, $\Cof_W\cap W=W^{\univ}$ $($\textit{set of co-universal $W$-equivalences in $\mathcal{M}$}$)$ \textit{is stable under filtering direct limits, and hence ``left stable'' in the sense of} (4.4.3).

\textbf{Corollary 4.4.4.} \textit{Under the conditions of proposition} 4.4.2, \textit{suppose moreover that simple amalgamated sums exist in $\mathcal{M}$, and let $C_0$ be a small subset of} $\Cof_W$ \textit{whose elements are accessible, and consider}
\[
C\;=\;\widetilde{C_0}\ ,\qquad TC\;\deq\;C\cap W\ .
\]
\textit{Then we have}
\[
C\;=\;\Cof_W\ ,\qquad TC\;=\;C\cap W^{\univ}\;\subset\;W^{\univ}\ ,
\]
\msp{page 41}
\textit{and $C$ and $TC$ are left $Q$-stable.}

Indeed, by theorem 1 of paragraph 1 (p. 10) and its corollary, $\widetilde{C_0}$ is the smallest subset of $\Fl(\mathcal{M})$ containing $C_0$ and left Q-stable. Since $\Cof_W$ is left Q-stable and contains $C_0$, we therefore have $C=\widetilde{C_0}\subset\Cof_W$. So $TC\deq C{\cap}W\subset\Cof_W{\cap}W\deq W^{\univ}$, so $TC=C\cap W^{\univ}$. Since $W^{\univ}$ is itself also left $Q$-stable (preceding corollary), the same is true of $C\cap W$, \quad \textit{q.e.d.}

\textbf{Corollary 4.4.5.} \textit{Suppose that every object of $\mathcal{M}$ is accessible. Under the conditions of the preceding corollary, setting}
\[
TF\;=\;C_{0*}\;=\;C_*\ ,\qquad F\;=\;(TC)_*\ ,
\]
\textit{the couple $(C,TF)$ is a Quillen couple. For the same to hold for the couple $(TC,F)$, it suffices that there exist a} small \textit{subset $TC_0$ of $C\cap W$, such that $\widetilde{TC_0}=TC$ $(\deq C\cap W)$.}

The first assertion is, as we saw earlier, theorem 1 of paragraph 1. The second follows from this same theorem. We can be more precise in this direction:

\textbf{Proposition 4.4.6.} \textit{Let $\mathcal{M}$ be a category stable under small filtering $\varinjlim$, under amalgamated sums, and all of whose objects are accessible, $W$ a localizer}
\msp{page 42}
\textit{mildly saturated in $\mathcal{M}$, stable under filtering $\varinjlim$, $C_0$ a} small \textit{subset of} $\Fl(\mathcal{M})$, \textit{such that}
\[
C_0\;\subset\;\Cof_W\ ,\qquad \underbrace{C_{0*}}_{(=TF)}\;\subset\;W\ ,
\]
\printedpage{28}
% head: A. Grothendieck : Dérivateurs, Ch. XV, Édition des Universités de Montpellier II et Paris VII
% foot: Version 5 novembre 2010
\textit{let us set}
\[
C\;=\;\widetilde{C_0}\qquad TC\;=\;C\cap W\ ,\qquad TF\;=\;C_*\ ,\qquad F\;=\;(TC)_*\ ,
\]
\textit{and let $TC_0$ be a small subset of $TC=C\cap W$. We then have}
\[
\left\{
\begin{aligned}
C &\;\subset\;\Cof_W\ , &\qquad TC &\;=\;C\cap W^{\univ}\quad(\textit{for the record})\\
\widetilde{TC_0} &\;\subset\;TC\ , & F_0 &\supset F\quad(\textit{the latter trivial})\ ,
\end{aligned}
\right.
\tag*{(4.4.6)}
\]
\textit{and the following conditions are equivalent:}
\begin{enumerate}[leftmargin=*]
\item[a)] $\widetilde{TC_0}=TC\quad(\textit{we have}\subset)$\,;
\item[a$'$)] $F_0=F\quad(\textit{we have}\supset)$\,;
\item[b)] $TF=F_0\cap W\quad(\textit{we have}\subset)$\footnote[14]{\textbf{N.B.} We see easily that $\boxed{TF=F\cap W,}$ but that does not suffice to have $TF=F_0\cap W$ for a suitable $TC_0$.}\,;
\item[c)] $W=TF\circ\widetilde{TC_0}\quad(\textit{we have}\supset)$\footnote[15]{\textbf{N.B.} We know that $\boxed{W=TF\circ TC,}$ which is in no way sufficient to imply $(c)$\,!}\,;
\end{enumerate}
\textit{and imply that $(W,C,F)$ is a closed Quillen triple.}

\textsc{Proof.} The inclusion $\widetilde{TC_0}\;\subset\;TC$ follows from the left Q-stability of $TC$ (cor. 4.4.4) and from the fact that $TC_0\subset TC$ is small (one then applies theorem 1 just as for $C\subset\Cof_W\ldots$). The equivalence of $(a)$, $(b)$, $(c)$ is none other than proposition 4.3 applied to $W$ and the two Quillen couples $(C,TF)$, $(\widetilde{TC_0},F_0)$ (taking into account theorem 1 of paragraph 1). That we then have a closed Quillen triple was seen in corollary 3 of proposition 4.3 (p. 35). It remains to prove the equivalence of $(a)$ and $(a')$. We have $(a)\Longrightarrow(a')$, since $(TC_0)_*=(\widetilde{TC_0})_*=F_0$. Conversely, $(a')$ implies $\widetilde{TC_0}=\widetilde{TC}$, hence the factorization theorem for $(TC,F)$ by virtue of the factorization theorem, corollary 1 (p. 11) applied to $TC_0\subset TC\subset\widetilde{TC_0}$ and to $F_0=F=(TC_0)_*$.

\msp{page 43}
Thus, the only problem, once $C_0$ (small) has been found satisfying $C_0\subset\Cof_W$, $C_{0*}\subset W$, is to find a \textit{small} subset $TC_0\subset TC$ which ``generates $TC\;=\;C\cap W$ '', \textit{i.e.} which satisfies the equivalent conditions $(a)$, $(a')$, $(b)$, $(c)$ of the preceding proposition. I have spent days trying to prove the existence of $TC_0$, assuming $\mathcal{M}$ equipped with a cardinal filtration, and $W$ an accessible subset of $\mathcal{M}$, without having succeeded. Of course, all that really matters to us is that $(TC,F)$ satisfy the factorization condition. But I see no way of proving it, except by ``generating'' $TC$ on the left by a \textit{small} subset $TC_0$.

\section*{5\quad Preliminary technical results on accessible subsets of $\mathrm{Ob}\,\mathcal{M}$, stable under filtering direct limits}

Yet I think I can arrive at something usable, by working systematically with $C$'s ($\subset\Cof_W$, left Q-closed, left Q-generated by a \textit{small} subset $C_0$ of $C$) which are stable under all filtering direct limits.

\end{document}
